% Chunk: Chapter 12, section 12.2, Cancellation of Divergences.
% Source: physical PDF pp. 528--539 (printed pp. 505--516), including the
% p. 528 continuation from section 12.1 and ending immediately before the
% section 12.3 heading on physical p. 539.
% Coverage: equations (12.2.1)--(12.2.26), all unnumbered displays, Figures
% 12.2--12.3, and five ordinary footnotes.
% Uncertainty: none. The assignment's nominal p. 538 endpoint was corrected
% against the scan: section 12.2 continues through most of physical p. 539.

the amplitude corresponding to any graph is that power-counting should
give $D<0$ not only for the complete multiple integral for the whole
amplitude, but also for any subintegration defined by holding any one
or more linear combinations of the loop momenta fixed. (The graphs
shown in Figures~\ref{fig:12.1}(b) and
\ref{fig:12.1}(c) fail this test because $D\geq0$ for the
subintegrations in which only the momenta for the loops within the dotted
squares are integrated.) We will not repeat the rather long proof here,
because it is well treated in earlier books~\hyperref[ref:12.3]{[3]}, and
in any case the method of proof has little to do with how we actually do
calculations. The next section will describe how this requirement is
fulfilled.

\subsection{Cancellation of Divergences}\label{sec:12.2}

Consider a Feynman diagram, or part of a Feynman diagram, with positive
superficial degree of divergence, $D\geq0$. The part of the momentum space
integral where all internal momenta go to infinity together will then
diverge, like $\int^\infty k^{D-1}\,dk$. If we differentiate $D+1$ times
with respect to any external momentum, we lower the net number of momentum
factors in the integrand by $D+1$,
\footnote{For instance, if an internal scalar field line carries a momentum
$k+p$, where $p$ is a linear combination of external four-momenta and $k$ is
a four-momentum variable of integration, then the derivative of the
propagator $[(k+p)^2+m^2]^{-1}$ with respect to $p^\mu$ gives
$-2(k_\mu+p_\mu)[(k+p)^2+m^2]^{-2}$, which goes as $k^{-3}$ rather than
$k^{-2}$ for $k\to\infty$.}
and hence render this part of the momentum space integral convergent. There
may still be divergences arising from subgraphs, like those in
Figures~\ref{fig:12.1}(b) and \ref{fig:12.1}(c); for the
moment we will ignore this possibility, returning to it later in this
section. Since differentiation $D+1$ times renders the integral finite, it
follows that the contribution of such a graph or subgraph can be written as
a polynomial of order $D$ in external momenta, with divergent coefficients,
plus a finite remainder.

To see how this works without irrelevant complications, consider the
logarithmically divergent one-dimensional integral
\[
  \mathcal{J}(q)=\int_0^\infty\frac{dk}{k+q}
\]
with $D=1-1=0$. Differentiating once gives
\[
  \mathcal{J}'(q)=-\int_0^\infty\frac{dk}{(k+q)^2}=-\frac1q,
\]
so
\[
  \mathcal{J}(q)=-\ln q+c.
\]
The constant $c$ is obviously divergent, but the rest of the integral is
perfectly finite. In exactly the same way, we can evaluate the $D=1$
integral
\[
  \int_0^\infty\frac{k\,dk}{k+q}=a+bq+q\ln q
\]
with divergent constants $a$ and $b$.

Now, a polynomial term in external momenta is just what would be produced
by adding suitable terms to the Lagrangian: if a graph with $E_f$ external
lines of type $f$ has degree of divergence $D\geq0$, then the ultraviolet
divergent polynomial is the same as would be produced by adding various
interactions $i$ with $n_{if}=E_f$ fields of type $f$ and $d_i\leq D$
derivatives. If there already are such interactions in the Lagrangian, then
the ultraviolet divergences simply add corrections to the coupling constants
of these interactions. Hence these infinities can be cancelled by including
suitable infinite terms in these coupling constants. All that we ever
measure is the sum of the bare coupling constant and the corresponding
coefficient from one of the divergent polynomials, so if we demand that the
sum equals the (presumably finite) measured value, then the bare coupling
must automatically contain an infinity that cancels the infinity from the
divergent integral over internal momenta. (One qualification: where the
divergence occurs in a graph or subgraph with just two external lines, which
appears as a radiative correction to a particle propagator, we must demand
not that some effective coupling constant equals its measured value, but
rather that the complete propagator has a pole at the same position and with
the same residue as for free particles.) In this way, all infinities are
absorbed into a redefinition of couplings constants, masses, and fields.

For this renormalization program to work, it is essential that the Lagrangian
include \emph{all} interactions that correspond to the ultraviolet divergent
parts of Feynman amplitudes. (There are exceptions to this rule in
supersymmetric theories~\hyperref[ref:12.4]{[4]}.) The interactions in the
Lagrangian are, of course, limited by various symmetry principles, such as
Lorentz invariance, gauge invariance, etc., but these constrain the
ultraviolet divergences in the same way. (It takes some work to prove that
non-Abelian gauge symmetries constrain infinities in the same way that they
constrain interactions. This will be shown in Volume II.) In the general
case, there are no other limitations on the ultraviolet divergences, so
\emph{the Lagrangian must include every possible term consistent with
symmetry principles}.

But there is an important class of theories with only a finite number of
interactions, where the renormalization program also works. These are the
so-called renormalizable theories, whose interactions all have $\Delta_i\geq0$.
Equation~\eqref{eq:12.1.8} then gives
\[
  D\leq4-\sum_f E_f(s_f+1),
\]
so divergent polynomials arise in only a limited number of Feynman graphs or
subgraphs: those with few enough external lines so that $D\geq0$. The
contribution of such divergent polynomials is just the same as would be
produced by replacing the divergent graph or subgraph with a single vertex
arising from a term in the Lagrangian with $E_f$ fields of type $f$ and
$0,1,\ldots,D$ derivatives. But, comparing with Eq.~\eqref{eq:12.1.9}, we see that
\emph{these are precisely the same as the interactions that satisfy the
renormalizability requirement} $\Delta_i\geq0$, or in other words,
\[
  0\leq d_i\leq4-\sum_f n_{if}(s_f+1).
\]

In order for all infinities to cancel in a renormalizable theory, it is
usually necessary that \emph{all} renormalizable interactions that are
allowed by symmetries must actually appear in the Lagrangian.
\footnote{In addition, interactions and mass terms that are not allowed by
global symmetries may appear in the Lagrangian, as long as they are
superrenormalizable, that is, with $\Delta_i>0$. This is because the presence
of a superrenormalizable coupling lowers the degree of divergence, so that
the symmetry breaking does not affect those divergences that are cancelled
by the strictly renormalizable couplings with $\Delta_i=0$. Note that it is
the \emph{bare} strictly renormalizable couplings that would exhibit the
symmetry; renormalized couplings that are defined in terms of mass-shell
matrix elements generally show the effect of symmetry breaking.}
For instance, if there is a scalar (or pseudoscalar) field $\phi$ and fermion
field $\psi$ with interactions $\bar\psi\psi\phi$ (or
$\bar\psi\gamma_5\psi\phi$) then we cannot exclude an interaction $\phi^4$;
otherwise there would be no counterterm to cancel the logarithmic divergence
arising from fermion loops with four attached scalar or pseudoscalar lines.

Let's see in more detail how the cancellation of infinities works in the
simplest version of quantum electrodynamics. Equation~\eqref{eq:12.1.11} shows that
the only graphs or subgraphs that could possibly yield divergent integrals
are the following:

\par\noindent\textbf{$E_e=2$, $E_\gamma=1$}\par
This is the electron--photon vertex $\Gamma_\mu^{*(\ell)}(p',p)$. (The
superscript $\ell$ indicates that this includes only contributions from
graphs with loops.) It has $D=0$, so its divergent part is
momentum-independent. Lorentz invariance then only allows this divergent
constant to be proportional to $\gamma_\mu$, so
\begin{equation}
  \Gamma_\mu^{*(\ell)}=L\gamma_\mu+\Gamma_\mu^{(f)},
  \tag{12.2.1}\label{eq:12.2.1}
\end{equation}
with $L$ a logarithmically divergent constant, and $\Gamma_\mu^{(f)}$ finite.
This does not uniquely define the constant $L$, since we can always move a
finite term $\delta L\gamma_\mu$ from $\Gamma_\mu^{(f)}$ to $L\gamma_\mu$. To
complete the definition, we may note that as shown in Section 9.7, the
mass-shell matrix element of $\Gamma_\mu^*(p,p)$ and hence of
$\Gamma_\mu^{(f)}(p,p)$ between mass-shell Dirac spinors is proportional to
the same matrix element of $\gamma^\mu$, so we may define $L$ by the
prescription that
\begin{equation}
  \bar u(p,\sigma')\Gamma_\mu^{(f)}(p,p)u(p,\sigma)=0
  \tag{12.2.2}\label{eq:12.2.2}
\end{equation}
for $p^2+m_e^2=0$.

\par\noindent\textbf{$E_e=2$, $E_\gamma=0$}\par
This is the electron self-energy insertion $\Sigma^*(p)$. It has $D=1$, so
its divergent part is linear in the momentum $p^\mu$ carried by the incoming
and outgoing fermion. Lorentz invariance (including parity conservation)
will only allow it to be a function of $\sl{p}$, so we may write the loop
contribution as
\begin{equation}
  \Sigma^{*(\ell)}(p)=A-(i\sl{p}+m)B+\Sigma^{(f)}(\sl{p}),
  \tag{12.2.3}\label{eq:12.2.3}
\end{equation}
where $A$ and $B$ are divergent constants, and $\Sigma^{(f)}$ is finite.
Again, this does not uniquely define the constants $A$ and $B$, because we
can always shift $\Sigma^{(f)}$ by a finite first-order polynomial in
$\sl{p}$. We will define $A$ and $B$ by the prescription that
\begin{equation}
  \Sigma^{(f)}=\frac{\partial\Sigma^{(f)}}{\partial\sl{p}}=0
  \quad\text{for}\quad i\sl{p}=-m.
  \tag{12.2.4}\label{eq:12.2.4}
\end{equation}
Actually, $B$ is not a new divergent constant. As long as we use a
regularization procedure that respects current conservation, $\Gamma_\mu$
and $\Sigma$ will be related by the Ward identity~\eqref{eq:10.4.27}
\[
  \Gamma^\mu(p,p)=\gamma^\mu+i\frac{\partial}{\partial p_\mu}\Sigma(p)
\]
and therefore
\begin{equation}
  L\gamma_\mu+\Gamma_\mu^{(f)}(p,p)
  =B\gamma_\mu+i\frac{\partial\Sigma^{(f)}(p)}{\partial p^\mu}.
  \tag{12.2.5}\label{eq:12.2.5}
\end{equation}
Taking the matrix element of this equation between $\bar u(p,\sigma')$ and
$u(p,\sigma)$ and using Eqs.~\eqref{eq:12.2.2} and \eqref{eq:12.2.4}, we find
\begin{equation}
  L=B.
  \tag{12.2.6}\label{eq:12.2.6}
\end{equation}

\par\noindent\textbf{$E_\gamma=2$, $E_e=0$}\par
This is the photon self-energy insertion $\Pi_{\mu\nu}^*(q)$. It has $D=2$,
so its divergent part is a second-order polynomial in $q$. Lorentz invariance
only allows $\Pi_{\mu\nu}^*$ to take the form of a linear combination of
$\eta_{\mu\nu}$ and $q_\mu q_\nu$ with coefficients depending only on $q^2$,
so the loop contributions take the form
\[
  \Pi_{\mu\nu}^{*(\ell)}(q)
  =C_1\eta_{\mu\nu}+C_2\eta_{\mu\nu}q^2+C_3q_\mu q_\nu
  +\text{finite terms},
\]
where $C_1$, $C_2$, and $C_3$ are divergent constants. As long as we use a
regularization technique that respects current conservation, we must have
\[
  q^\mu\Pi_{\mu\nu}^{*(\ell)}(q)=0.
\]
The same must then be true for the divergent terms, so
$C_1q_\nu+(C_2+C_3)q^2q_\nu$ must be finite for all $q$. It follows that
$C_1$ and $C_2+C_3$ must be finite, and can therefore be lumped into the
finite part of $\Pi_{\mu\nu}^{*(\ell)}(q)$. Thus
\begin{equation}
  \Pi_{\mu\nu}^{*(\ell)}(q)
  =(\eta_{\mu\nu}q^2-q_\mu q_\nu)\bigl(C+\pi(q^2)\bigr),
  \tag{12.2.7}\label{eq:12.2.7}
\end{equation}
where $\pi(q^2)$ is finite and $C$ is the sole remaining divergence in
$\Pi_{\mu\nu}$. To pin down the definition of $C$, we may move any finite
constant $\pi(0)$ into $C$, so that
\begin{equation}
  \pi(0)=0.
  \tag{12.2.8}\label{eq:12.2.8}
\end{equation}

\par\noindent\textbf{$E_\gamma=4$, $E_e=0$}\par
This is the amplitude $M_{\mu\nu\rho\sigma}$ for scattering of light by
light. It has $D=0$, so using Lorentz invariance and Bose statistics, it may
be written (there is no non-loop contribution)
\[
  M_{\mu\nu\rho\sigma}
  =K(\eta_{\mu\nu}\eta_{\rho\sigma}
     +\eta_{\mu\rho}\eta_{\nu\sigma}
     +\eta_{\mu\sigma}\eta_{\nu\rho})+\text{finite terms}
\]
with $K$ a potentially divergent constant. However, current conservation
gives
\[
  q^\mu M_{\mu\nu\rho\sigma}=0
\]
and so $K(q_\nu\eta_{\rho\sigma}+q_\rho\eta_{\nu\sigma}
+q_\sigma\eta_{\nu\rho})$ is finite. In order for this to be true for
$q\neq0$, $K$ must itself be finite. This is a nice example of the role of
symmetry principles in the renormalization program: if $K$ had turned out
to be infinite it could not be removed by renormalization of the coupling
constant for an interaction $(A_\mu A^\mu)^2$, because no such interaction
is allowed by gauge invariance, but $K$ is finite because of current
conservation conditions that are imposed by gauge invariance.

\par\noindent\textbf{$E_\gamma=1$, $E_e=0$ and $E_e=1$, $E_\gamma=0,1,2$}\par
These have $D=3$ and $D=\tfrac52,\tfrac32$, and $\tfrac12$, respectively,
but Lorentz invariance makes all such graphs vanish.

\par\noindent\textbf{$E_\gamma=3$, $E_e=0$}\par
This has $D=1$, but vanishes because of charge-conjugation invariance.

The reader will perhaps have noticed that the independent divergent
constants $A$, $B$, $C$ are in one-to-one correspondence with the independent
parameters $Z_2$, $Z_3$, and $\delta m$ in the counterterm
part~\eqref{eq:11.1.9} of
the Lagrangian for quantum electrodynamics. These counterterms make a direct
contribution $Z_2\delta m-(Z_2-1)(i\sl{p}+m)$ to $\Sigma^*(p)$. The
requirement that the position and residue of the one-particle pole be the
same as in the free-field propagator means we must choose $Z_2$ and
$\delta m$ so that the total $\Sigma^*(p)$ satisfies Eq.~\eqref{eq:12.2.4},
i.e.,
\begin{align}
  Z_2\delta m&=-A,
  \tag{12.2.9}\label{eq:12.2.9}\\
  Z_2-1&=-B,
  \tag{12.2.10}\label{eq:12.2.10}
\end{align}
so that the complete electron self-energy insertion is just the finite
function $\Sigma^{(f)}(p)$:
\begin{equation}
  \Sigma(p)=\Sigma^{(f)}(p).
  \tag{12.2.11}\label{eq:12.2.11}
\end{equation}
Also, $\mathcal{L}_2$ makes a direct contribution to $\Gamma_\mu$ equal to
$(Z_2-1)\gamma_\mu$. Using Eq.~\eqref{eq:12.2.6}, we see that the full
vertex is
\begin{equation}
  \Gamma_\mu=\gamma_\mu+(Z_2-1)\gamma_\mu+\Gamma_\mu^{*(\ell)}
  =\gamma_\mu+\Gamma_\mu^{(f)}.
  \tag{12.2.12}\label{eq:12.2.12}
\end{equation}
This is not only finite, but satisfies the condition
\begin{equation}
  \bar u(p,\sigma')\Gamma_\mu(p,p)u(p,\sigma)
  =\bar u(p,\sigma')\gamma_\mu u(p,\sigma),
  \tag{12.2.13}\label{eq:12.2.13}
\end{equation}
as can also be seen from Eqs.~\eqref{eq:10.6.13} and
\eqref{eq:10.6.14}. Finally,
$\mathcal{L}_2$ makes a contribution
$-(Z_3-1)(q^2\eta_{\mu\nu}-q_\mu q_\nu)$ to $\Pi_{\mu\nu}^*(q)$. In order
that the photon propagator should have a pole with the same residue as for
free fields we need the coefficient of $q^2\eta_{\mu\nu}-q_\mu q_\nu$ in
the total $\Pi_{\mu\nu}(q)$ to vanish, so
\begin{equation}
  Z_3=1+C
  \tag{12.2.14}\label{eq:12.2.14}
\end{equation}
and the photon propagator is then finite:
\begin{equation}
  \Pi_{\mu\nu}(q)=(\eta_{\mu\nu}q^2-q_\mu q_\nu)\pi(q^2).
  \tag{12.2.15}\label{eq:12.2.15}
\end{equation}

So far, we have only checked that the divergences, arising from the region
of momentum space in which all internal momenta are large (and with generic
ratios), are polynomials in external momenta that are cancelled by suitable
counterterms. Such graphs are called \emph{superficially convergent}. Before
we conclude that all ultraviolet divergences actually are removed by
renormalization, we need to consider the ultraviolet divergences arising in
higher-order graphs when some subset of the momentum space integration
variables rather than all of them go to infinity. For instance, in quantum
electrodynamics the superficial divergences in subintegrations come from
subgraphs that are either photon self-energy parts $\Pi^*$, or electron
self-energy parts $\Sigma^*$, or electron--electron--photon vertices
$\Gamma^\mu$.

% Figure 12.2: overlapping divergences in the photon self-energy.
% Source: physical PDF p. 534 (printed p. 511).
\begingroup
\renewcommand{\thefigure}{12.\arabic{figure}}
\setcounter{figure}{1}
\begin{figure}[tbp]
\centering
\begin{tikzpicture}[
  electron/.style={line width=0.7pt,
    postaction={decorate},decoration={markings,
      mark=at position 0.50 with {\arrow{Stealth[length=4.6pt,width=3.5pt]}}}},
  electronq/.style={line width=0.7pt,
    postaction={decorate},decoration={markings,
      mark=at position 0.26 with {\arrow{Stealth[length=4.6pt,width=3.5pt]}},
      mark=at position 0.74 with {\arrow{Stealth[length=4.6pt,width=3.5pt]}}}},
  photon/.style={line width=0.65pt,decorate,
    decoration={snake,amplitude=1.0pt,segment length=4.8pt}},
  callout/.style={-{Stealth[length=4.6pt,width=3.5pt]},line width=0.65pt},
  every node/.style={font=\small}
]
  % Overlapping two-loop photon self-energy.
  \begin{scope}[xshift=-3.25cm,yshift=1.35cm]
    \draw[photon] (-3.00,0) -- (-1.45,0);
    \draw[photon] (1.45,0) -- (3.00,0);
    \draw[electronq] (-1.45,0) arc[start angle=180,end angle=0,x radius=1.45cm,y radius=0.88cm];
    \draw[electronq] (1.45,0) arc[start angle=0,end angle=-180,x radius=1.45cm,y radius=0.88cm];
    \draw[photon] (0,0.88) -- (0,-0.88);
  \end{scope}
  % Counterterm at the left photon-electron vertex.
  \begin{scope}[xshift=3.25cm,yshift=1.35cm]
    \draw[photon] (-3.00,0) -- (-1.45,0);
    \draw[photon] (1.45,0) -- (3.00,0);
    \draw[electron] (-1.45,0) arc[start angle=180,end angle=0,x radius=1.45cm,y radius=0.88cm];
    \draw[electron] (1.45,0) arc[start angle=0,end angle=-180,x radius=1.45cm,y radius=0.88cm];
    \draw[line width=0.8pt] (-1.57,-0.12)--(-1.33,0.12);
    \draw[line width=0.8pt] (-1.57,0.12)--(-1.33,-0.12);
    \draw[callout] (-2.05,-1.15) -- (-1.62,-0.48);
    \node at (-1.55,-1.43) {$(Z_2-1)_2$};
  \end{scope}
  % Counterterm at the right photon-electron vertex.
  \begin{scope}[xshift=-3.25cm,yshift=-1.45cm]
    \draw[photon] (-3.00,0) -- (-1.45,0);
    \draw[photon] (1.45,0) -- (3.00,0);
    \draw[electron] (-1.45,0) arc[start angle=180,end angle=0,x radius=1.45cm,y radius=0.88cm];
    \draw[electron] (1.45,0) arc[start angle=0,end angle=-180,x radius=1.45cm,y radius=0.88cm];
    \draw[line width=0.8pt] (1.33,-0.12)--(1.57,0.12);
    \draw[line width=0.8pt] (1.33,0.12)--(1.57,-0.12);
    \draw[callout] (1.95,-1.12) -- (1.62,-0.45);
    \node at (1.72,-1.42) {$(Z_2-1)_2$};
  \end{scope}
  % Photon-field counterterm.
  \begin{scope}[xshift=3.25cm,yshift=-1.45cm]
    \draw[photon] (-3.00,0) -- (3.00,0);
    \draw[line width=0.8pt] (-0.12,-0.13)--(0.12,0.13);
    \draw[line width=0.8pt] (-0.12,0.13)--(0.12,-0.13);
    \draw[callout] (-0.55,-1.10) -- (-0.12,-0.42);
    \node at (0,-1.42) {$(Z_3-1)_{\mathrm{overlap}}$};
  \end{scope}
\end{tikzpicture}
\caption{Some fourth-order graphs for the photon self-energy in quantum electrodynamics that involve overlapping divergences. Lines carrying arrows are electrons; wavy lines photons. The crosses mark the contribution of counterterms.}
\label{fig:weinberg-12-2}
\end{figure}
\endgroup



The problem with such divergences is that they cannot be removed by
differentiation with respect to external momenta; we are left with terms
where the derivatives act only on internal lines in the parts of the graph
which are \emph{not} in the divergent subgraphs, and therefore do not reduce
the degree of divergence of these subgraphs. As mentioned in the previous
section, a graph or sum of graphs is actually convergent only if it and all
its subintegrations are superficially convergent, in the sense of counting
powers of momentum. But wherever such a divergent subgraph appears, it
always comes accompanied with an infinite counterterm. In electrodynamics,
these are the terms in Eq.~\eqref{eq:11.1.9}: a term
$-(Z_3-1)(q^2\eta_{\mu\nu}-q_\mu q_\nu)$ for each $\Pi_{\mu\nu}^*(q)$; a
term $Z_2\delta m-(Z_2-1)(i\sl{p}+m)$ for each $\Sigma^*(p)$; and a term
$(Z_2-1)\gamma^\mu$ for each $\Gamma^\mu$. Just as for the graph as a whole,
these counterterms cancel the infinities from the divergent subgraphs.

Unfortunately, there is a flaw in this simple argument---the possibility of
overlapping divergences. That is, it is possible that two divergent subgraphs
may share an internal line, so that we cannot regard them as independent
divergent integrals. In quantum electrodynamics this happens only when two
electron--electron--photon vertices overlap inside a photon or an electron
self-energy insertion,
\footnote{The sharing of a line in two self-energy insertions or in a
self-energy insertion and a vertex part would not leave enough external
lines to attach such a subgraph to the rest of the diagram. Historically,
the Ward identity~\eqref{eq:10.4.26} was used to by-pass the problem of overlapping
divergences in the electron self-energy, by expressing the electron
self-energy in terms of the vertex function, where overlapping divergences
do not occur. This approach will not be followed here, as it is unnecessary,
and in any case does not solve the problem for the self-energy of the photon
or other neutral particles.}
as shown in Figures~\ref{fig:weinberg-12-2} and \ref{fig:weinberg-12-3}.

% Figure 12.3: overlapping divergences in the electron self-energy.
% Source: physical PDF p. 535 (printed p. 512).
\begingroup
\renewcommand{\thefigure}{12.\arabic{figure}}
\setcounter{figure}{2}
\begin{figure}[tbp]
\centering
\begin{tikzpicture}[
  electron/.style={line width=0.7pt,
    postaction={decorate},decoration={markings,
      mark=at position 0.52 with {\arrow{Stealth[length=4.6pt,width=3.5pt]}}}},
  photon/.style={line width=0.65pt,decorate,
    decoration={snake,amplitude=1.0pt,segment length=4.8pt}},
  callout/.style={-{Stealth[length=4.6pt,width=3.5pt]},line width=0.65pt},
  every node/.style={font=\small}
]
  % Overlapping two-loop electron self-energy.
  \begin{scope}[xshift=-3.25cm,yshift=1.35cm]
    \draw[line width=0.7pt] (-3.00,0) -- (-1.45,0);
    \draw[line width=0.7pt] (1.45,0) -- (3.00,0);
    \draw[electron] (-1.45,0) .. controls (-1.18,0.68) and (-0.52,0.88) .. (0,0.88);
    \draw[photon] (-1.45,0) .. controls (-1.18,-0.68) and (-0.52,-0.88) .. (0,-0.88);
    \draw[line width=0.7pt] (0,0.88) -- (0,-0.88);
    \draw[photon] (0,0.88) .. controls (0.52,0.88) and (1.18,0.68) .. (1.45,0);
    \draw[electron] (0,-0.88) .. controls (0.52,-0.88) and (1.18,-0.68) .. (1.45,0);
  \end{scope}
  % Vertex counterterm at the left end of the loop.
  \begin{scope}[xshift=3.25cm,yshift=1.35cm]
    \draw[line width=0.7pt] (-3.00,0) -- (-1.45,0);
    \draw[line width=0.7pt] (1.45,0) -- (3.00,0);
    \draw[photon] (-1.45,0) arc[start angle=180,end angle=0,x radius=1.45cm,y radius=0.88cm];
    \draw[electron] (-1.45,0) arc[start angle=180,end angle=360,x radius=1.45cm,y radius=0.88cm];
    \draw[line width=0.8pt] (-1.57,-0.12)--(-1.33,0.12);
    \draw[line width=0.8pt] (-1.57,0.12)--(-1.33,-0.12);
    \draw[callout] (-2.05,-1.15) -- (-1.62,-0.48);
    \node at (-1.55,-1.43) {$(Z_2-1)_2$};
  \end{scope}
  % Vertex counterterm at the right end of the loop.
  \begin{scope}[xshift=-3.25cm,yshift=-1.45cm]
    \draw[line width=0.7pt] (-3.00,0) -- (-1.45,0);
    \draw[line width=0.7pt] (1.45,0) -- (3.00,0);
    \draw[electron] (-1.45,0) arc[start angle=180,end angle=0,x radius=1.45cm,y radius=0.88cm];
    \draw[photon] (1.45,0) arc[start angle=0,end angle=-180,x radius=1.45cm,y radius=0.88cm];
    \draw[line width=0.8pt] (1.33,-0.12)--(1.57,0.12);
    \draw[line width=0.8pt] (1.33,0.12)--(1.57,-0.12);
    \draw[callout] (1.95,-1.12) -- (1.62,-0.45);
    \node at (1.72,-1.42) {$(Z_2-1)_2$};
  \end{scope}
  % Electron-field counterterm.
  \begin{scope}[xshift=3.25cm,yshift=-1.45cm]
    \draw[line width=0.7pt] (-3.00,0) -- (3.00,0);
    \draw[line width=0.8pt] (-0.12,-0.13)--(0.12,0.13);
    \draw[line width=0.8pt] (-0.12,0.13)--(0.12,-0.13);
    \draw[callout] (-0.55,-1.10) -- (-0.12,-0.42);
    \node at (0,-1.42) {$(Z_3-1)_{\mathrm{overlap}}$};
  \end{scope}
\end{tikzpicture}
\caption{Some fourth-order graphs for the electron self-energy in quantum electrodynamics that involve overlapping divergences. Wavy lines are photons, other lines are electrons. The crosses mark the contribution of counterterms.}
\label{fig:weinberg-12-3}
\end{figure}
\endgroup


A complete treatment of renormalization that takes account of overlapping
divergences should include a prescription for eliminating superficial
ultraviolet divergences, not only in the overall integration but in all
subintegrations as well, together with a proof that this prescription is (at
least formally) implemented by renormalization of masses, fields, and
coupling constants. The theorem of Ref.~\hyperref[ref:12.2]{[2]} then
ensures that all Green's functions of renormalized fields are finite when
expressed in terms of renormalized masses and couplings. The first proof that
the renormalization of fields, couplings, and masses renders the whole
integration and all its subintegrations superficially divergent was offered
by Salam~\hyperref[ref:12.5]{[5]}. A more specific prescription for
eliminating ultraviolet divergences was given by Bogoliubov and
Parasiuk~\hyperref[ref:12.6]{[6]}, and corrected by
Hepp~\hyperref[ref:12.7]{[7]}, and shown by them to be equivalent to a
renormalization of fields, masses, and coupling constants. Finally,
Zimmermann~\hyperref[ref:12.8]{[8]} proved that this prescription does
eliminate all superficial divergences in the whole integration and all its
subintegrations, and used the theorem of Ref.~\hyperref[ref:12.2]{[2]} to
conclude from this that the renormalized Feynman momentum-space integrals
are convergent.

Briefly, the `BPHZ' prescription for eliminating superficial divergences
requires that we consider all possible ways (called `forests') of surrounding
a whole graph and/or its subgraphs with boxes that may be nested within each
other but do not overlap. (An example is given below.) For each forest we
define a subtraction term by replacing the integrand for any subgraph of
superficial divergence $D$ within a box (starting with the innermost boxes and
working outwards) with the first $D+1$ terms of its Taylor series expansion
in the momenta flowing into or out of that box.
\footnote{As described here, this prescription applies to unrenormalizable as
well as renormalizable theories. In renormalizable theories it implies that
there is no subtraction unless the box contains one of the limited number of
graphs corresponding to renormalizable terms in the Lagrangian.}
The subtracted Feynman amplitude is given by the original graph, minus all
these subtraction terms, including the subtraction term for a forest
consisting of a single box surrounding the whole graph.

It is fairly easy to see that the subtracted Feynman amplitude that is
calculated in this way is the same as would be obtained by replacing all
fields, coupling constants, and masses in the original Lagrangian by their
renormalized counterparts. The difference between this procedure and the
sort of renormalization we carried out in Chapter 11 is that the renormalized
fields, coupling constants, and masses are defined in terms of amplitudes at
an unconventional renormalization point, where all four-momenta vanish. (In
this respect, the one-dimensional divergent integrals discussed at the start
of this section provide an elementary example of the BPHZ method of
separating divergent terms.) But there is nothing special about this
renormalization point; once a Feynman amplitude is made convergent by
expressing it in terms of these unconventional renormalized quantities, it
can be rewritten in terms of conventionally renormalized fields, couplings,
and masses without introducing new infinities.

It is not necessary to use the BPHZ subtraction prescription in practice.
Replacing fields, masses, and couplings with their renormalized counterparts
(defined using any convenient renormalization points) automatically provides
counterterms that cancel all infinities. Instead of proving that the BPHZ
subtraction prescription really does make all integrals converge, we shall
just look at one example that shows how renormalization works, even in the
presence of overlapping divergences.

Consider the fourth-order contribution to the photon self-energy insertion
$\Pi_{\mu\nu}^*(q)$ shown in Figure~\ref{fig:weinberg-12-2}. (The forests here
consist of the whole integral over $p$ and $p'$, the subintegration over $p$
alone, and the subintegration over $p'$ alone.) Including the corresponding
counterterms for the vertex parts and photon field renormalization, this has
the value
\begin{align}
 [\Pi_{\mu\nu}^*(q)]_{\mathrm{overlap}}
 ={}&-\frac{e^4}{(2\pi)^8}\int d^4p\int d^4p'\,
 \frac{1}{(p-p')^2-i\epsilon}
 \notag\\
 &\quad\cdot\operatorname{Tr}\!\left\{
 S(p')\gamma_\nu S(p'+q)\gamma^\rho
 S(p+q)\gamma_\mu S(p)\gamma_\rho\right\}
 \notag\\
 &-2(Z_2-1)_2\frac{i e^2}{(2\pi)^4}\int d^4p\,
 \operatorname{Tr}\!\left\{\gamma_\nu S(p+q)\gamma_\mu S(p)\right\}
 \notag\\
 &-(Z_3-1)_{\mathrm{overlap}}
 (q^2\eta_{\mu\nu}-q_\mu q_\nu),
 \tag{12.2.16}\label{eq:12.2.16}
\end{align}
where $S(p)\equiv[-i\sl{p}+m]/[p^2+m^2-i\epsilon]$;
$(Z_2-1)_2$ is the term in $Z_2-1$ of second order in $e$; and
$(Z_3-1)_{\mathrm{overlap}}$ is a logarithmically divergent constant of
fourth order in $e$ that cancels the terms in
$[\Pi_{\mu\nu}^*(q)]_{\mathrm{overlap}}$ of second order in $q^\mu$. The
factor 2 in the second term arises because there is a renormalization
counterterm $Z_2-1$ for each of the two vertices in the second-order photon
self-energy. Note, however, that the first term here can be thought of
\emph{either} as the insertion of a vertex correction given by the
$p'$-integral into a photon self-energy given by the $p$-integral, or as the
insertion of a vertex correction given by the $p$-integral into a photon
self-energy given by the $p'$-integral, but \emph{not} as the insertion of
two independent vertex corrections, because there is only one photon
propagator.

To see how to handle the infinities in Eq.~\eqref{eq:12.2.16}, note that
\begin{equation}
  \bigl[(Z_2-1)_2+R_2\bigr]\gamma_\mu
  =\frac{i e^2}{(2\pi)^4}\int\frac{d^4p'}{p'^2-i\epsilon}\,
  \gamma_\rho S(p')\gamma_\mu S(p')\gamma^\rho,
  \tag{12.2.17}\label{eq:12.2.17}
\end{equation}
where $R_2$ is a finite remainder. (Lorentz invariance tells us that the
integral on the right is proportional to $\gamma_\mu$. The difference
between this integral and $(Z_2-1)_2\gamma_\mu$ equals the complete
renormalized electron--electron--photon vertex to second order in $e$ at
zero electron and photon momenta, and is therefore finite.) This allows us
to rewrite Eq.~\eqref{eq:12.2.16} in the form
\begin{align}
 [\Pi_{\mu\nu}^*(q)]_{\mathrm{overlap}}
 ={}&-\frac{e^4}{(2\pi)^8}\int d^4p\int d^4p'
 \Biggl[
 \frac{1}{(p-p')^2-i\epsilon}
 \notag\\
 &\qquad\cdot
 \operatorname{Tr}\!\left\{S(p')\gamma_\nu S(p'+q)\gamma^\rho
 S(p+q)\gamma_\mu S(p)\gamma_\rho\right\}
 \notag\\
 &\qquad-\frac{1}{p'^2-i\epsilon}
 \operatorname{Tr}\!\left\{S(p')\gamma_\nu S(p')\gamma^\rho
 S(p+q)\gamma_\mu S(p)\gamma_\rho\right\}
 \notag\\
 &\qquad-\frac{1}{p^2-i\epsilon}
 \operatorname{Tr}\!\left\{S(p')\gamma_\nu S(p'+q)\gamma^\rho
 S(p)\gamma_\mu S(p)\gamma_\rho\right\}
 \Biggr]
 \notag\\
 &-2R_2\frac{i e^2}{(2\pi)^4}\int d^4p\,
 \operatorname{Tr}\!\left\{\gamma_\nu S(p+q)\gamma_\mu S(p)\right\}
 \notag\\
 &-(Z_3-1)_{\mathrm{overlap}}
 (q^2\eta_{\mu\nu}-q_\mu q_\nu).
 \tag{12.2.18}\label{eq:12.2.18}
\end{align}
First consider integration over $p'$ alone. Each of the first two terms is
logarithmically divergent, but their difference is finite. The third term is
also logarithmically divergent (with a gauge-invariant regulator), but the
divergence in this term (unlike the first two terms) takes the form of a
second-order polynomial in $q$, with the remainder finite. This remaining
divergence is cancelled by the term
$-(Z_3-1)(q^2\eta^{\mu\nu}-q^\mu q^\nu)$ that cancels all second-order terms
in the expansion of $\Pi_{\mu\nu}^*(q)$. So the $p'$-subintegration gives a
finite result. The symmetry of Eq.~\eqref{eq:12.2.16} shows that in exactly
the same way the $p$ subintegration also gives a finite result. Generic
subintegrations over $p$ and $p'$ with $ap+bp'$ held fixed (where $a$ and
$b$ are arbitrary non-zero constants) are manifestly convergent, and the
integration over $p$ and $p'$ together is made finite by the counterterm
$-(Z_3-1)(q^2\eta^{\mu\nu}-q^\mu q^\nu)$. Thus Eq.~\eqref{eq:12.2.18} and any of its
subintegrations satisfy the power-counting requirements for convergence,
and therefore according to the theorem~\hyperref[ref:12.2]{[2]} quoted in
the previous section, the whole expression actually converges.

\begin{center}
  $*\quad *\quad *$
\end{center}

In electrodynamics there is a natural definition of renormalized couplings
as well as renormalized masses and fields. This is not always the case. For
example, consider the theory of a single real scalar field $\phi(x)$ with
Lagrangian density
\begin{equation}
  \mathcal{L}=-\frac12\partial_\mu\phi\partial^\mu\phi
  -\frac12m^2\phi^2-\frac1{24}g\phi^4.
  \tag{12.2.19}\label{eq:12.2.19}
\end{equation}
To one-loop order, the S-matrix for scalar--scalar scattering is given by the
Feynman rules as
\begin{equation}
  S(q_1q_2\to q'_1q'_2)
  =\frac{-i(2\pi)^4\delta^4(p'_1+p'_2-p_1-p_2)}
  {(2\pi)^6(16E'_1E'_2E_1E_2)^{1/2}}
  F(q_1q_2\to q'_1q'_2),
  \tag{12.2.20}\label{eq:12.2.20}
\end{equation}
where
\begin{align}
 -i(2\pi)^4F(q_1q_2\to q'_1q'_2)
 ={}&-i(2\pi)^4g
 +\frac12\bigl[-i(2\pi)^4g\bigr]^2
 \left[\frac{-i}{(2\pi)^4}\right]^2
 \notag\\
 &\cdot\int d^4k\Biggl[
 \frac{1}{[(q_1+k)^2+m^2-i\epsilon][(q_2-k)^2+m^2-i\epsilon]}
 \notag\\
 &\hspace{4.0cm}+(q_2\to-q'_1)+(q_2\to-q'_2)\Biggr],
 \tag{12.2.21}\label{eq:12.2.21}
\end{align}
and $q_1,q_2$ and $q'_1,q'_2$ are the incoming and outgoing four-momenta.
Combining denominators and rotating the $k^0$-integration contour as usual,
this is
\begin{align}
 F={}&g-\frac{g^2}{16\pi^2}\int_0^\infty k^3\,dk\int_0^1dx
 \Bigl\{[k^2+m^2-sx(1-x)]^{-2}
 \notag\\
 &\qquad+[k^2+m^2-tx(1-x)]^{-2}
       +[k^2+m^2-ux(1-x)]^{-2}\Bigr\},
 \tag{12.2.22}\label{eq:12.2.22}
\end{align}
where $s,t$, and $u$ are the Mandelstam variables
\begin{equation}
  s=-(p_1+p_2)^2,\qquad
  t=-(p_1-p'_1)^2,\qquad
  u=-(p_1-p'_2)^2,
  \tag{12.2.23}\label{eq:12.2.23}
\end{equation}
related by $s+t+u=4m^2$; also, $x$ is the Feynman parameter introduced in
combining denominators. With an ultraviolet cutoff at $k=\Lambda$, this gives
the result (for $\Lambda\gg m$)
\begin{align}
 F=g-\frac{g^2}{32\pi^2}\int_0^1dx\Biggl\{
 &\ln\!\left(\frac{\Lambda^2}{m^2-sx(1-x)}\right)
 +\ln\!\left(\frac{\Lambda^2}{m^2-tx(1-x)}\right)
 \notag\\
 &+\ln\!\left(\frac{\Lambda^2}{m^2-ux(1-x)}\right)-3\Biggr\}.
 \tag{12.2.24}\label{eq:12.2.24}
\end{align}
We can define the renormalized coupling $g_R$ as the value of $F$ at any
point $s,t,u$ we like, provided we stay in the region where $F$ is real. For
instance, suppose that in order to maintain the symmetry among the scalars,
we choose to renormalize at the off-mass-shell point
\footnote{Going back over the derivation of Eq.~\eqref{eq:12.2.25}, one may
check that in this derivation we have not used the conditions
$p_1^2=p_2^2=p_1'^2=p_2'^2=-m^2$, so Eq.~\eqref{eq:12.2.24} is valid
whatever we take for the external line masses.}
$p_1^2=p_2^2=p_1'^2=p_2'^2=\mu^2$, $s=t=u=-4\mu^2/3$. Defining the
renormalized coupling $g_R$ as the value of $F$ at this point, we have
\begin{equation}
  g=g_R+\frac{3g_R^2}{32\pi^2}
  \left[\ln\!\left(\frac{\Lambda^2}{\mu^2}\right)-1
  -\int_0^1dx\,\ln\!\left(\frac{4x(1-x)}3+\frac{m^2}{\mu^2}\right)\right]
  +\cdots.
  \tag{12.2.25}\label{eq:12.2.25}
\end{equation}
The cutoff dependence then cancels in Eq.~\eqref{eq:12.2.24} to order
$g_R^2$, leaving a finite formula for $F$ in terms of $g_R$:
\begin{align}
 F=g_R-\frac{g_R^2}{32\pi^2}\int_0^1dx\Biggl\{
 &\ln\!\left(\frac{m^2+4x(1-x)\mu^2/3}{m^2-sx(1-x)}\right)
 +\ln\!\left(\frac{m^2+4x(1-x)\mu^2/3}{m^2-tx(1-x)}\right)
 \notag\\
 &+\ln\!\left(\frac{m^2+4x(1-x)\mu^2/3}{m^2-ux(1-x)}\right)
 \Biggr\}+\cdots.
 \tag{12.2.26}\label{eq:12.2.26}
\end{align}
Here $\mu^2$ may be taken to be any real quantity greater than $-3m^2$, in
which range $g_R$ is real. The explicit $\mu$-dependence in
Eq.~\eqref{eq:12.2.26} is, of course, cancelled by the $\mu$-dependence of
the renormalized coupling. This freedom to change the renormalization
prescription (which of course exists also in electrodynamics and other
realistic theories) will be of great importance to us when we come to the
renormalization group method in Volume II.
